\subsection{G. The three graded flatness questions}

Let $K$ be any commutative ring, and let $A,B,C$ be nonnegatively graded
commutative $K$-algebras with $A_0=B_0=C_0=K$.  Let
$\psi:A\to B$ and $\varphi:B\to C$ be graded $K$-algebra maps.
The literal ``good conditions'' on cote 23, scans 096--098, say that
$\varphi$ has a graded $K$-linear section $u:C\to B$, and that
\[
 v:A\otimes_K C\longrightarrow B,\qquad
 a\otimes c\longmapsto\psi(a)u(c)
\]
is bijective.  We regard its source as a graded $A$-module.  The map $v$
is $A$-linear, whether or not $u$ is multiplicative.  We do not assume
finite generation of any homogeneous component.

\begin{theorem}[Proved: exact flatness classification]\label{SA23-G-1}
Under these good conditions:
\begin{align*}
 B\text{ is flat over }A
 &\quad\Longleftrightarrow\quad C\text{ is flat over }K,\\
 B\text{ is flat over }K
 &\quad\Longleftrightarrow\quad
 A\text{ and }C\text{ are flat over }K,\\
 C\text{ is flat over }B
 &\quad\Longleftrightarrow\quad A=K.
\end{align*}
The first two equivalences also hold degree by degree for the assertion
that all homogeneous components are projective over $K$ in the second
line.
\end{theorem}
\begin{proof}
The map $v$ identifies $B$ with the induced $A$-module $A\otimes_K C$.
Flatness of $C$ over $K$ thus implies flatness of $B$ over $A$.
Conversely, apply the base-change functor $K\otimes_A-$, where
$A\to K=A/A_+$ is the augmentation.  The resulting $K$-module is
$C$, and base change preserves flatness.

If $A$ and $C$ are $K$-flat, their tensor product is $K$-flat.
Conversely, $A=A\otimes C_0$ and $C=A_0\otimes C$ are direct
summands of $A\otimes C$ as $K$-modules, using the homogeneous
decompositions.  A direct summand of a flat module is flat.  The same
argument in degree $n$, using
$B_n\simeq\bigoplus_{i+j=n}A_i\otimes_K C_j$, proves the projective
version: the tensor product of two projective modules over a commutative
ring is projective, and the displayed sum is finite.

For the third equivalence put $I=\ker\varphi$.  It is a homogeneous
ideal with $I_0=0$.  The exact sequence $0\to I\to B\to C\to0$
gives $\operatorname{Tor}_1^B(C,C)\simeq I/I^2$.  Thus $B$-flatness
of $C$ implies $I=I^2$.  A positively graded ideal with this equality
is zero: in degree one its square is zero, and induction on the degree
then gives $I_n=0$ for every $n$.  Consequently $\varphi$ is an
isomorphism and $u=\varphi^{-1}$.  For $a\in A_n$, $n>0$, put
$c=\varphi\psi(a)\in C_n$.  The element
$a\otimes1-1\otimes c$ lies in the kernel of $v$, so it is zero.
Its two terms belong to the distinct direct summands $(n,0)$ and
$(0,n)$ of $(A\otimes C)_n$.  Therefore $a=0$.  We have proved
$A_+=0$.  If $A=K$, on the other hand, $v=u$ is an isomorphism;
$\varphi=u^{-1}$, so $C$ is free of rank one over $B$.
\end{proof}

\begin{example}[Refuted: the first implication on scan 098]\label{SA23-G-2}
Let $K=\mathbf F_2$, $A=B=\mathbf F_2[e]/(e^2)$ with $\deg e=1$,
$\psi=\operatorname{id}$, and $C=\mathbf F_2$ with
$\varphi(e)=0$.  Take $u$ to be the inclusion of constants.  Then
$v:A\otimes_K C\to B$ is an isomorphism and every graded piece is
finite free over $K$.  Thus even the ``very good'' conditions hold and
$A$ is $K$-flat.  But $C$ is not $B$-flat: tensoring
$(e)\hookrightarrow B$ with $C$ gives the zero map from the
one-dimensional space $(e)/(e)^2$ to $C$.  This verifies every hypothesis
and disproves $A/K$ flat $\Rightarrow C/B$ flat.  The four-element ring
and this tensor calculation are exhaustively checked in
\texttt{code/algebra\_finite\_checks.py}.
\end{example}

\begin{proposition}[Proved with an explicit hypothesis: section independence]\label{SA23-G-3}
Assume additionally that $\varphi\psi$ kills $A_+$.  Then the following
are equivalent: $\ker\varphi=\psi(A_+)B$; and $v$ is surjective.
If $v$ is bijective for one graded linear section, it is bijective for
every such section.  If every $A_n,B_n,C_n$ is finite projective and
$v$ is surjective, equality of the fiberwise Hilbert series
$P_B=P_A P_C$ implies that $v$ is bijective.
\end{proposition}
\begin{proof}
If the kernel equality holds, an element of $B_n$ is the sum of the
lift $u(c)$ of its image in $C_n$ and terms $\psi(a_i)b_i$ with
$\deg a_i>0$ and $\deg b_i<n$.  Induction on $n$ makes $v_n$
surjective.  Conversely, in an expression for an element of
$\ker\varphi$ obtained from $v$, the component from $A_0\otimes C_n$
must be zero after applying $\varphi$, and all other components lie in
$\psi(A_+)B$.

Let $u'$ be another section.  The difference $u'(c)-u(c)$ lies in the
kernel.  Under $v^{-1}$ its terms have positive $A$-degree, and hence
strictly smaller $C$-degree than $c$.  In total degree $n$ the
endomorphism $v^{-1}v'$ is $1+N$ with $N^{n+1}=0$.
Its inverse is $1-N+\cdots+(-N)^n$.  This proves independence.
Finally, a surjective map of finite projective modules having the same
rank at every prime is an isomorphism: split the surjection, and its
finite projective kernel has rank zero everywhere, hence is zero.
Apply this separately in every degree.
\end{proof}

\begin{remark}[Refuted literal readings of scans 096--097]\label{SA23-G-4}
The supplied transcription does not explicitly state
$\varphi\psi(A_+)=0$.  Without it, the claimed equivalence with the
kernel condition is false: take $K=\mathbf F_2$,
$A=B=C=K[e]/(e^2)$, and both maps and the section to be identities.
Then $v$ is surjective while
$\ker\varphi=0\ne(e)=\psi(A_+)B$.
Likewise, the Hilbert-series inequality does not follow from the
splitting condition alone: $A=C=\mathbf F_2$, $B=\mathbf F_2[e]/(e^2)$, with augmentation
$B\to C$, contradicts it in degree one.  It follows from the splitting
\emph{and} the kernel condition (or directly from surjectivity of $v$).
Theorem~\ref{SA23-G-1}, which answers the three questions, is valid under
the literal good conditions and does not need this additional hypothesis.
Both finite examples are checked in the attached program.
These are mathematical alternatives to be distinguished; until we
inspect the original scan we do not assert which omission is in the
handwriting.
\end{remark}
